\chapter{The SIR Model, Continued}

\section{The threshold phenomenon}

Say the initial proportion of infecteds is $I(0) \ne 0$. Will the
infection \emph{invade}? I.e.~will $I$ increase? Equivalently, does
$\dot I>0$ hold?

\[
\dot I = \beta S I - \gamma I = I(\beta S - \gamma).
\]
So  $\dot{I} > 0$ if 
\[
\beta S - \gamma  > 0.
\]
We must have $\beta S - \gamma > 0$, or equivalently
\[
S > \frac{\gamma}{\beta}.
\]
So if $S<\gamma/\beta$ the infection cannot invade. The implication
for vaccination strategies and herd immunity is immediate.

\subsection{The basic reproduction number}
What is $\beta/\gamma$? Recall $1/\gamma$ is the average time to
recovery, i.e.\ the disease duration and $\beta$ is infections per
individual per time unit (in a fully susceptible population). Hence
$R_0$ is the expected number of secondary infections caused by a
single infected individual introduced into an otherwise fully
susceptible population. We denote it
\[
R_0 = \frac{\beta}{\gamma}
\]


\section{Demographics}
We want to include births and deaths. Individuals are born into the
susceptible class $S$ at rate $\mu$ (per unit of the whole population)
and die at the same per-capita rate $\mu$ from every class, so
population size is still fixed. The new equations are
\begin{align*}
\dot S &= \mu - \beta S I - \mu S, \\
\dot I &= \beta S I - \gamma I -\mu I, \\
\dot R &= \gamma I - \mu R.
\end{align*}
As before, $\dot S + \dot I + \dot R = 0$, so $S+I+R=1$ stays
constant.

\subsection{$R_0$ for an SIR model with demographics}
With demography, an infected individual leaves the infected class
either by recovering (rate $\gamma$) or by dying (rate $\mu$),
independently. Recall that \emph{minimum} of two independent
exponentials is exponential with the sum of rates, so the time spent
infected is $\text{Exp}(\gamma+\mu)$.

Using the same interpretation as before, $R_0$ is (contacts during
infection) $\times$ (transmission probability per contact) with the
mean infectious period $1/(\gamma+\mu)$:
\[
R_0 = \frac{\beta}{\gamma + \mu}.
\]

\section{Equilibria}
We find equilibrium states by equating the derivatives to zero.

\subsection{Equilibrium analysis for basic SIR model}

Recall the basic SIR model:
\begin{align*}
\dot S &= -\beta S I, \\
\dot I &= \beta S I - \gamma I, \\
\dot R &= \gamma I.
\end{align*}
Set all three derivatives to zero. $\dot R = \gamma I = 0$ forces
$I=0$ (since $\gamma>0$). Then  $\dot S = \dot{I} = 0$.

So the equilibria are $(S,0,R)$ for $S+R=1$ (and, of course, $0\le
S\le 1$). This is a \emph{disease-free equilibrium} (DFE).

\subsection{Equilibrium analysis for SIR model with demographics}

Recall the model with demography:
\begin{align*}
\dot S &= \mu - \beta S I - \mu S, \\
\dot I &= \beta S I - \gamma I - \mu I, \\
\dot R &= \gamma I - \mu R.
\end{align*}
Set all three derivatives to zero. From $\dot I=0$:
$I[\beta S-(\gamma+\mu)]=0$, so either $I=0$ or
$S=\dfrac{\gamma+\mu}{\beta}=\dfrac{1}{R_0}$.

\paragraph{Disease-free equilibrium (DFE).} Then $\dot S=\mu(1-S)=0$
forces $S=1$. Other compartments are empty.

\paragraph{Endemic equilibrium (EE).} Take $\dot I=0$. If $I = 0$ we are at DFE. So assume $I > 0$. Then
\[
\dot{I} = 0 \Rightarrow \beta S^* = \gamma + \mu \Rightarrow S^* = \frac{\gamma +\mu}{\beta} = \frac{1}{R_0}.  
\]
Then
\[
I^* = \frac{\mu}{\gamma+\mu}\left(1-\frac{1}{R_0}\right)
\]
and
\[
R^* = \frac{\gamma}{\gamma+\mu}\left(1-\frac{1}{R_0}\right).
\]
So
\[
(S^*,I^*,R^*) = \left(\frac{1}{R_0},\ \ \frac{\mu}{\gamma+\mu}\Big(1-\frac{1}{R_0}\Big),\ \ \frac{\gamma}{\gamma+\mu}\Big(1-\frac{1}{R_0}\Big)\right).
\]
Note that this equilibrium is feasible (i.e.\ $I^*>0$) exactly when
$R_0>1$.

\section{Stability analysis}

\subsection{Linearization}

For a real function of one variable, near a point $x^*$:
\[
f(x^*+\delta x) \approx f(x^*) + \frac{df}{dx}(x^*)\,\delta x.
\]
In two dimensions a function takes two inputs and produces two
outputs $(x,y)\mapsto (f(x,y),g(x,y))$. So:

\begin{align*}
f(x^*+\delta x,\,y^*+\delta y) &\approx f(x^*,y^*) + \frac{\partial f}{\partial x}\delta x + \frac{\partial f}{\partial y}\delta y \\
g(x^*+\delta x,\,y^*+\delta y) &\approx g(x^*,y^*) + \frac{\partial g}{\partial x}\delta x + \frac{\partial g}{\partial y}\delta y.
\end{align*}
Collecting:
\[
\begin{pmatrix} f(x^*+\delta x,y^*+\delta y) \\ g(x^*+\delta x,y^*+\delta y) \end{pmatrix}
\approx
\begin{pmatrix} f(x^*,y^*) \\ g(x^*,y^*) \end{pmatrix}
+
\begin{pmatrix} \partial f/\partial x & \partial f/\partial y \\ \partial g/\partial x & \partial g/\partial y \end{pmatrix}
\begin{pmatrix} \delta x \\ \delta y \end{pmatrix}.
\]

\begin{definition}
The \emph{Jacobian} of $(f,g)$ at $(x^*,y^*)$ is
\[
J(x^*,y^*) = \begin{pmatrix} \partial f/\partial x & \partial f/\partial y \\ \partial g/\partial x & \partial g/\partial y \end{pmatrix},
\]
\end{definition}
So:
\[
\begin{pmatrix} f(x^*+\delta x,y^*+\delta y) \\ g(x^*+\delta x,y^*+\delta y) \end{pmatrix}
\approx
\begin{pmatrix} f(x^*,y^*) \\ g(x^*,y^*) \end{pmatrix}
+
J(x^*, y^*) \begin{pmatrix} \delta x \\ \delta y \end{pmatrix}.
\]

Now, our goal is to linearize a dynamical system 
\[
\begin{pmatrix} \dot x \\ \dot y \end{pmatrix} = \begin{pmatrix} f(x,y) \\ g(x,y) \end{pmatrix}
\]
using the Jacobian. We linearize at an equilibrium point, $f(x^*,y^*)
= g(x^*,y^*) = 0$ and let $x=x^*+\delta x$, $y=y^*+\delta y$. Then
\begin{align*}
  \begin{pmatrix}
    \dot x \\ \dot y
  \end{pmatrix} &= \begin{pmatrix} f(x^*,y^*) \\ g(x^*,y^*) \end{pmatrix} \\
\end{align*}

Since $x^*,y^*$ are constant, $\dot x=\dot{\delta x}$ and $\dot
y=\dot{\delta y}$, so by the approximation above and using $f(x^*,y^*)
= g(x^*,y^*) = 0$:
\[
\begin{pmatrix} \dot{\delta x} \\ \dot{\delta y} \end{pmatrix}
\approx
J(x^*,y^*)
\begin{pmatrix} \delta x \\ \delta y \end{pmatrix}.
\]
Near an equilibrium, the dynamics of a small perturbation are
therefore approximately \emph{linear}, governed entirely by the
Jacobian.

\paragraph{Why is this useful?}
Stability of an equilibrium is ``stable'' determined by its
eigenvalues. Roughly, all eigenvalues have to be
negative\footnote{have negative real part.} for an equilibrium to be
stable.

\subsection{Linearization of SIR model}
Here the variables are $S$ and $I$:
\begin{align*}
\dot S &= f(S,I) := \mu(1-S)-\beta SI \\
\dot I &= g(S,I) := \beta SI - (\gamma+\mu)I\\ 
\end{align*}
The Jacobian is
\[
J(S,I) = \begin{pmatrix} -\mu-\beta I & -\beta S \\ \beta I & \beta S-(\gamma+\mu) \end{pmatrix}.
\]


\subsubsection{Linearization around the disease-free equilibrium (DFE)}
In the disease-free equilibrium $I^*=0$ and $S^*=1$:
\[
J(S^*, I^*) = \begin{pmatrix} -\mu & -\beta \\ 0 & \beta-(\gamma+\mu) \end{pmatrix}
\]
Eigenvalues are:
\[
\lambda_1 = -\mu, \qquad \lambda_2 = \beta-(\gamma+\mu) = (\gamma+\mu)(R_0-1).
\]
$\lambda_1<0$ always; $\lambda_2<0 \iff R_0<1$. So the disease-free
equilibrium is stable iff $R_0<1$.

\subsubsection{Linearization around the endemic equilibrium (EE)}
In the endemic equilibrium,
\[
(S^*,I^*,R^*) =
\left(\frac{1}{R_0},\ \ \frac{\mu}{\gamma+\mu}\Big(1-\frac{1}{R_0}\Big),\ \ \frac{\gamma}{\gamma+\mu}\Big(1-\frac{1}{R_0}\Big)\right).
 \]
$\beta S^*=\gamma+\mu$ and $\beta
I^*=\mu(R_0-1)$ and the Jacobian becomes
\[
J(S^*,I^*) = \begin{pmatrix} -\mu R_0 & -(\gamma+\mu) \\ \mu(R_0-1) & 0 \end{pmatrix},
\qquad
\operatorname{tr}J(S^*,I^*) = -\mu R_0, \qquad \det J(S^*,I^*) = \mu(\gamma+\mu)(R_0-1).
\]
For a $2\times 2$ matrix, both eigenvalues have negative real part iff
$\operatorname{tr}<0$ and $\det>0$. Here $\operatorname{tr}J^*<0$
always, and $\det J^*>0 \iff R_0>1$. So whenever the endemic
equilibrium exists ($R_0>1$), it is automatically stable.

\paragraph{Summary.}
$R_0<1 \Rightarrow$ disease-free equilibrium is stable. $R_0>1 \Rightarrow$ endemic equilibrium is stable.
